% ============================================================
%  F. C. Sibbern, On the Concept, Nature, and Essence of Philosophy.
%    A Presentation of the Propaedeutic of Philosophy.
%  Copenhagen: At the Author's Own Expense (Printed by Director
%    Jens Hostrup Schultz), 1843.
%  TRANSLATION (English) — SCAFFOLD.
%  Source of truth for translation: transcription.tex (Danish, COMPLETE + proofed,
%    §§1--21 + Fortale + Rettelser; verified compile 59 pp., 0/0).
%  Method: ../../../TRANSLATION-PLAYBOOK.md. Fill each "% [text to be added: ...]"
%    marker batch by batch (~10 printed pp.), mirroring transcription 1:1.
%  Same policy as the companion volume *Om Erkjendelse og Granskning*:
%    - Register: lightly modernized (faithful meaning, Sibbern's long periods untangled).
%    - Key terms: on first use in a section, gloss the Danish in brackets, e.g.
%      "fundamental cognition [Grunderkjendelse]". Maintain GLOSSARY.md.
%    - Layout: English-only, section-keyed; mirrors the transcription structure 1:1.
%  Page-offset (Danish scan, for reference): printed page = PDF page - 11.
%  Book-specific: emphasis in this book is BOLD Fraktur (Fett), already encoded as
%    \emph{} in the transcription -> keep as \emph{} (renders italic). Latin already in
%    \textit{}; German block-quotes are plain (Fraktur in the source); no real Greek.
%  See TRANSLATION-NOTES.md for state, decisions, and the section->page map.
% ============================================================
\documentclass[12pt,a4paper]{book}
\usepackage[utf8]{inputenc}
\usepackage[T1]{fontenc}
\usepackage{libertinus}
\usepackage{libertinust1math}
\usepackage{textalpha}   % real Greek glyphs if any arise (only transliterated names here)
\usepackage[english]{babel}
\usepackage{enumitem}    % for \arabic*)-style lists (playbook: easy to forget)
\usepackage[protrusion=true,expansion=true]{microtype}
\usepackage{setspace}
\usepackage[margin=1.2in]{geometry}
\usepackage{fancyhdr}
\setlength{\headheight}{28pt}
\addtolength{\topmargin}{-13.5pt}
\usepackage{hyperref}

\hypersetup{
  pdftitle={On the Concept, Nature, and Essence of Philosophy},
  pdfauthor={F. C. Sibbern},
  hidelinks
}

\pagestyle{fancy}
\fancyhf{}
\fancyhead[LE,RO]{\thepage}
\fancyhead[CE]{\textit{\leftmark}}
\fancyhead[CO]{\textit{\rightmark}}

\setstretch{1.3}

% Group-heading helper (the centered bold rubrics of the Indhold / body)
\newcommand{\grouphead}[1]{%
  \bigskip\begin{center}\large\textbf{#1}\end{center}\smallskip}

% TITLE — rendering to confirm with Hans (see TRANSLATION-NOTES).
\title{\textbf{On}\\[6pt]
  \Huge\textbf{the Concept, Nature,}\\[6pt]
  \LARGE\textbf{and Essence of Philosophy.}\\[14pt]
  \normalsize A Presentation of\\[4pt]
  \Large\textbf{the Propaedeutic of Philosophy.}\\[16pt]
  \normalsize By\\[6pt]
  \large Dr.\ Frederik Christian Sibbern,\\[4pt]
  \normalsize Professor of Philosophy.}
\author{}
\date{Copenhagen, 1843.\\
  At the Author's Own Expense\\
  printed by Director Jens Hostrup Schultz,\\
  Royal and University Printer.}

\begin{document}
\maketitle
\thispagestyle{empty}
\cleardoublepage

%% ============================================================
%%  PREFACE  (Fortale, PDF 8--9). Transcription L78--111.
%% ============================================================
\chapter*{Preface.}
\markboth{Preface}{Preface}

\begin{center}\rule{0.28\linewidth}{0.4pt}\end{center}

Had I not, up to now, let every writing of mine that has come forward
independently in the literature be accompanied by a preface---not exactly
because I have expressly reckoned such a thing among the proper appointments of
a book, but because each time it has felt natural to me---I should this time
have been inclined to dispense with any preface at all, since I had nothing to
say that could be its subject. For the truth is that I would rather have
refrained from expressly declaring that what several of my readers will perhaps
expect to find stated here I mean to leave to them to settle for themselves. I
mean this: in what relation the present writing stands to my book
\textit{Om Erkjendelse og Granskning} [On Cognition and Inquiry], published at
the close of the year 1821. If I am to declare myself on the matter---for how I
intend to use both writings in my lectures on the propaedeutic of philosophy
[Philosophiens Propædeutik] does not belong here to explain---then, so far as
the public is concerned, I would say nothing more than this: I could wish that,
in its relation to the present book, the reader would regard that earlier one
as though it were by another author, one whom I had drawn upon and to whom I had
therefore often referred, while yet having a different object before me than he.
For in that earlier work the object is science [Videnskab] and scientific study
in general, and it is offered as an introduction to academic study at large.
Here, by contrast, I have had philosophy in particular before me, and have
therefore had more to enter into with respect to \emph{this} than had come
forward in that other work. ---Finally, I ask that among the errata [Rettelser],
those marked \textbf{NB} in particular may be heeded.

\bigskip
\noindent\textbf{Copenhagen,} the 28th of October 1843.
\nopagebreak
\begin{flushright}\textbf{Sibbern.}\end{flushright}

\begin{center}\rule{0.18\linewidth}{0.4pt}\end{center}

\mainmatter

%% ============================================================
%%  PRELIMINARY REMARKS  (§§ 1--2, printed pp. 1--6)
%% ============================================================
\grouphead{Preliminary Remarks.}

\begin{center}\large\S~1.\end{center}
\markboth{\S~1}{\S~1}
\begin{center}What is to be undertaken here, and where we set out from.\end{center}

It is the concept, nature, and essence of philosophy that we shall here seek to
clarify for ourselves. We cannot, then, confine ourselves to the question: What
is meant by philosophy? but have to do with the far more telling question: What
is philosophy? Yet we must of course, straightaway and from the very
first, take account of the word's meaning, attending to what is called
philosophy---and has through the course of history been so called, and has
asserted itself under that name. For although the concept we are to seek to give
is to be developed and determined as reflection on the matter itself demands,
still the question---whether what is thereby set up corresponds to what one
calls philosophy, or whether perhaps, with all the validity it may have in and
for itself, it is something quite other than what one has in view by the word
philosophy---can only be settled by taking into consideration the meaning with
which this name passes and holds good. By the name, the object to be spoken of
is pointed out, and one must from the first take care not to end up speaking of
some other object instead of this one. It is most correct, then, to set out from
a given representation [Forestilling] of what philosophy is---a given meaning of
the word.

But the matter itself, thus pointed to by the name and thereby lifted out for
us, we must then surely consider entirely as it wills to be seen, and must
pursue it through its whole consequence, so that the concept, further on, cannot
be bound to, nor in its whole carrying-through be conditioned by, that given
representation.

In this way we do indeed come to consider philosophy in philosophy's own light,
and thus come to lay our own philosophy at the foundation; for we are now to
philosophize our way to the cognition [Erkjendelse] we seek, and in all
philosophizing philosophy at once asserts itself. But this circularity is a
consequence of the whole position of the matter; for philosophy is but one
member within the world-totality [Verdenstotalitet] which philosophy itself---%
according to the concept of it that we must presuppose as given---is to clarify
for us.

{\small\emph{Note~1.} There naturally arises here the question whether the
concept of philosophy ought not to be treated in two quite different places:
first, in the considerations introductory to philosophy, where the philosophical
undertaking cannot of course be dealt with without laying at its foundation a
concept of what philosophy is; ---and second, in the very midst of philosophy's
system, once the speculative fundamental cognition [Grunderkjendelse] has been so
far attained and developed that philosophy---stepping forward as a co-engaging
member in the spiritual life unfolding in consciousness---comes to be considered
from the speculative---more precisely a psycho-biological---standpoint. One
readily sees that, although philosophy would lead us back to the first
fundamental moments and beginning (\textit{initia}) of all things, it
nevertheless itself belongs in the very midst of the world's course, where it
can first arise in the higher life of consciousness, or in the midst of the
sphere of spirituality, and must then show itself as an essential principal
moment in the totality of spiritual life. The question of philosophy's proper
essence [Væsen] belongs, then, there where poetry, morality, and religiosity too
are to constitute themselves for us in their essential belonging-together with
the totality of spiritual life. ---\par}

{\small The question then naturally arises whether we ought not to distinguish,
from this spirit-philosophical and properly speculative consideration of
philosophy's essence, another, preliminary one belonging to a mere introduction
to philosophy---which we might then call the merely propaedeutic, and which would
at the same time be called the merely explicative. This one would merely pursue
the given concept of philosophy through all its consequence and thereafter
determine and expound it, whereas the other would let philosophy itself
constitute itself for us and thereby show itself in its proper essence. This last
would then be the one which, if the very expressive designation the philosophy of
philosophy is to be adopted, would have to be furnished with that name.\par}

{\small Yet to such a merely propaedeutic exposition of the concept of philosophy
we neither can nor will here confine ourselves, nor keep it separated from the
speculative; for it is philosophy itself---not the concept that has formed about
it---that we have before us here, and the clarification of philosophy, if it is
to be complete, cannot rest at that merely propaedeutic level.\par}

{\small In the meantime we must surely seek constantly to remain aware how far
our concept, as it more and more determines and develops itself for us, remains
within the domain [Gebeet] of the merely propaedeutic---so that it can be set up
as valid with respect to the various philosophical parties (Platonists,
Skeptics, Kantians, Lockeans, and so forth), and can thus be laid at the
foundation in the history of philosophy---and how far, on the other hand, it goes
further, and so is already developed according to the dictate of a definite
philosophical fundamental view [Grundanskuelse].\par}

{\small\emph{Note~2.} One will readily see that it would be altogether
inappropriate to give the first and fundamental part of philosophy itself---which
Hegel has called Logic---the name of the philosophy of philosophy, as some have
done. Hegel, too, is very far from having been able to give any occasion for
this, since with him the philosophy of philosophy---though this name is not found
in him---appears at the close of the third and last part of his whole
system.\par}

{\small\emph{Note~3.} Concerning the expression \emph{nature and essence} and its
relation to the expression \emph{concept}, be it noted that what is here most
immediately and first sought is the concept of philosophy and its exact
determination, but that this concept is then to unfold itself for us, or to step
out into the whole sum-total [Indbegreb] of its consequence, or into its whole
content---which is what we have in view by the expression \emph{nature and
essence}. (I refer here to my \textit{Logik}, \S~35.)\par}

\begin{center}\large\S~2.\end{center}
\markboth{\S~2}{\S~2}
\begin{center}A survey of the whole course of what follows.\end{center}

If, then, we ask first of all: What does one understand by philosophy?---whereby
we are only to indicate provisionally just so much that we may know with
certainty what it is that is to be treated, but must also at once come to
indicate the first points to which we have to direct our attention---then we note
that the received concept implies that philosophy is to aim at a certain
fundamental cognition [Grunderkjendelse] serving for a deeper clarification of
the whole of existence [Tilværelse], but that it is to form itself through
thinking and reflection, and thus to become a cognition that is to be aware of
its own inner groundedness and able to give an account of it. In this way, then,
we come to let our consideration take the following course:

First and foremost we take into consideration that 1) philosophy is a whole body
of cognizing [Indbegreb af Erkjenden], and we therefore consider the nature and
essence of cognition in general.

Next we observe that, since philosophy aims at a certain deeper fundamental
cognition that looks toward the last or innermost fundamental relations
[Grundforhold]---yet such a cognition shows itself as philosophy only when it
comes forward as sprung from rational cognition [Fornufterkjendelse], and
especially only when investigation [Efterforskning] or inquiry [Randsagelse] has
thereby been set in motion---we are hereby led to consider the endeavor of
cognition in general in its transition from its first principal stage, the stage
of first, immediate cognizing, to its second principal stage, the mediated but at
the same time higher stage of cognizing called forth by the spirit of inquiry
[Randsagelsesaand]---the stage at which cognition becomes a cognizing according
to insight, and at the same time---since here testing must take place, an account
must be rendered, grounds come under consideration, and groundedness is brought
to light and made good---a cognition in force of grounds [i Kraft af Grunde];
which then lets us see philosophy as that which 2) is a whole body of intelligent
and rational cognizing.

To this attaches, further, the observation that, since this whole body must form
a well-connected totality [Heelhed], philosophy in its carrying-through and
execution must become 3) a scientific cognition---whereby we are brought to
consider science [Videnskab] and scientific character [Videnskabelighed] in
general.

But since---according to the received concept---not all scientific cognizing is
philosophy, then, by asking ourselves what sort of science it is in which
philosophy consists, we come to become aware that the endeavor to press through
to the grounds cannot but lead to the forming of 4) a \emph{science} that in a
certain \emph{all-comprehending manner} aims at ultimate grounds and an ultimate
fundamental cognition, and that it is this above all that one has in view when
one speaks of philosophy; and here, when the question is now raised whether it is
the ultimate grounds and fundamental relations of all cognition, or of the whole
of existence, that are thereby meant, it will appear that---although a twofold
principal task [Hovedopgave] for philosophy has in this respect formed, and had
to form---it is nevertheless in the end only 5) the ultimate grounds and
fundamental relations of \emph{existence}, or as one is wont to say, of \emph{all
things}, in the finding-out and cognition of which the all-comprehending
fundamental cognition is to form, wherein philosophy consists.

This determination of the concept of philosophy, handed down from ancient times,
from which we are thus brought to set out, shows us philosophy as that which
comes forward over against, and in a certain opposition to, the sum-total of all
that cognizing whereby everything is given us and presents itself to us---%
everything whose ultimate grounds and fundamental relations are now to be in
question, and which we are accordingly prompted and called upon to gather for
ourselves into a single overview.

In its pursuit through the various principal respects or principal directions, we
shall next see the concept of philosophy lead to a whole cycle of conceptual
determinations, which can be developed into a whole cycle of definitions.

Thereupon we shall proceed to consider philosophy in its execution. But when we
now first take into consideration what ultimate grounds are, and how matters
stand with our pressing toward them, the question of philosophy's foundation
[Grundlag] and starting-point will prove to have to precede the question of the
execution of philosophy's work---to which, then, once we have thereupon fully
answered for ourselves the question what philosophy is, a consideration of
philosophy's significance in the totality of existence will naturally be able to
succeed.

With this, the course we shall take is fully indicated.

{\small\emph{Note.} It must be noted here at once---what will be more closely
treated further on in its place---that alongside philosophy proper, of which alone
there is question here, there is a whole body of sciences which, without being
philosophy proper, are nevertheless also called philosophical, but which I, to
distinguish them from the former (see further on, \S~11), give the name of
explicative philosophy, and to which I here pay no regard whatever.\par}

\begin{center}\rule{0.18\linewidth}{0.4pt}\end{center}

%% ============================================================
%%  INTRODUCTORY CONSIDERATIONS. ON COGNITION IN GENERAL. (§§ 3--7, pp. 7--24)
%% ============================================================
\grouphead{Introductory Considerations.\\[2pt]\normalsize On Cognition in General.}

\begin{center}\large\S~3.\end{center}
\markboth{\S~3}{\S~3}
\begin{center}Philosophy is a) cognition. ---The fundamental moments of all cognition.\end{center}

All cognizing aims at an \emph{is}, whereby something, in asserting itself in us,
is posited by us as an objective [det Objective] to which we must ascribe
validity, and which, insofar as it comes forward for us in this its validity, we
call the true---which is then to be taken entirely as it gives itself to us, that
is: shows itself to be.

Here, then, we see an opposition [Modsætning] between a self-conscious
subjectivity and an objectivity subsisting independently of it---one which,
moreover, comes forward for us not only from without but also from within. But
cognition sets in only when this opposition is so overcome that what is comes
forward for the subjectivity too exactly as it is, so that it comes to be for
us---that is, is posited by us to be---just as it is in itself; whereby we come,
in the very act of living in our own individual subsistence, at the same time to
live in the fullness of objectivity, indeed of the true.

The mediating and unifying element between the true and the subjectivity---which
latter shows itself at the same time as an individuality---is, on the side of the
first, the determining necessity with which it asserts itself, and on the side of
the subjectivity it is \emph{partly} the latter's becoming-determined and the
coming-to-consciousness that thereby sets in, which may be called a receiving
(\textit{scimus, qvia accepimus}), \emph{partly} the individual life's own doing
that answers to that becoming-determined. ---But this last is again partly one
that aims merely at a directing of the mind [Sindshenvendelse], partly one that
aims at inner subjective re-presentation and objectification. In the first we see
a subjective reaction and response, called forth by the awakening power of the
determining influence, that aims merely to make the reception deeper and
fuller---both through an accompanying devotion and through a supervening
recollection---as a care that the true may, in all its points and relations, come
to its rightful full power in us, so that in each and all we may be determined
wholly in accordance with the true's own bearing. In the subjective
re-presentation and objectification, by contrast, we see a productivity springing
from the individual life itself, which has its inner ground and mode of activity
within itself, and now, as a producing doing, brings forth and constitutes within
the subjective self-consciousness the inner presentation, in and through which the
true comes to be taken up into the subjective life. (And here, properly, it is
first said: \textit{scimus, qvia facimus}. ---For the rest I refer here to my work
\textit{Om Erkjend.\ og Granskn.}, \S\S~1--3, as well as to my \textit{Psychologie}.)
But what makes it that in our own productivity we nonetheless have, not something
merely subjective, but---by means of and through this merely subjective---the true
itself, is the power with which the true so determines and governs the subjective
life that creates and forms the representation, that it thereby at the same time
presents itself as the valid for the self-conscious beholding.

{\small\emph{Note~1.} When we here speak of an ``is,'' we see here not only an
opposition between an ``is'' and an ``is-not,'' but also an opposition between an
``is'' and a ``shows itself'' or ``seems.'' Hereupon one may at once ask whether
this ``seems'' too must not be said to be, and conversely, whether an ``is'' can
ever be anything other than a ``shows itself.'' Here it would be a matter of
asking what that ``is'' really means; and when we then observe that we can in
general have to do only with a being insofar as it asserts itself, but that now
every seeming too asserts itself, then the question is whether there is here a
differing self-assertion, and on what such difference rests---whereby a difference
between a merely momentary and subjective and a more constant and objective
self-assertion may perhaps be assumed, and within the latter there may again be
distinguished a) a merely sporadic, b) a merely microcosmic, c) a macrocosmic
validity, and only the last is to be acknowledged as the properly true one in the
last instance. But here the consideration would arise that all apparent, all
sporadic, all microcosmic self-assertion must in the end be grounded in the
macrocosmic order of things, so that here it is to be noted that all that other
self-assertion is seen in its true light---that is, in its proper being---only by
being seen in its emergence from, and reducibility to, the macrocosmic order of
things. ---But all the considerations that might thus be called forth by
reflection on that ``is'' do not belong here, but to the fundamental
considerations of speculative philosophy itself. ---This much is certain: that
throughout all our striving for cognition, the opposition between an ``is'' and a
``seems''---an ``in itself'' and a ``not in itself''---and along with it the
opposition between an objective and a subjective, so asserts itself that in the
whole course of our cognizing, as in the whole of the rest of our life bound up
with it, we cannot withdraw ourselves from it. And to this we hold ourselves
here, where we have merely to expound the nature of cognition as it gives itself
to us, and where we take without further ado the existence of a cognizing, with
all the moments therein, as given---just as we take the existence of a philosophy
as given---without entering here into the question of the proper reality of all
cognition and all philosophy, which it belongs to speculative philosophy itself
to decide. ---The concept of cognition we have here given has plainly---this is
still to be noted---validity even for one who would deny all our cognizing any
true reality; for such a one too must nonetheless have a concept of cognition that
he lays at the foundation, and it can only be the one given here; and one who
would ascribe to all our cognizing only a certain subjective validity must still
admit that cognition rests on the agreement of our representing with something
that takes on the character of an objective.\par}

{\small\emph{Note~2.} As an ``is'' forms itself, intelligence at once intervenes,
and a so-called category---or better, an a priori fundamental determination---%
asserts itself. Herewith cognition too at once comes forward as that which issues
in a judging, whereby a something (as subject) comes forward as that which is
something definite (the fundamental determination of whatness), and this
what-it-is comes forward as an other, which is nonetheless, by the judging itself,
made one with the former (as predicate). This fundamental nature of judging is
expressed in the so-called principle of identity.\par}

{\small\emph{Note~3.} The question arising here---how that productivity, having
its inner ground and mode of activity, and hence also the regulative principle for
it, within itself, nonetheless comes to present a real objective---is of
speculative nature, and finds for us its answer in the Idea of a
\textit{harmonia originaria} between individuality and totality. The determining
influence of the objective is indeed only an eliciting one, since it is the
individual life itself that, by its own activity within the individuality's circle
of consciousness, brings about that whereby the objective is to become the
cognized. But just as the whole individuality first comes forth at all through the
total course of things, and becomes, in its whole existence, as that course brings
it about, so too its whole organization---both in inner, psychic, and in outer
respects---is constituted in such a way, answering to the totality, that every
influence coming from the totality, whether it come from without or from the inner
universal, calls forth its own determinate corresponding activity in the
individual organization, or in the microcosm, and this so that it now, according
to the latter's own fundamental nature, leads precisely to such a certain
determinate production that just that determinate objective is thereby represented
by us. So that, although we constantly behold objects in our own soul, we
nonetheless behold the actual objects as they are in themselves, because our inner
representations become images of them, or counter-images to them, inasmuch as our
inner life of consciousness---according to that inner harmony of its with the
totality, when it (be it noted) is in its normal course---must each time bring
forth, out of its own inner working, precisely that nerve-action, with
corresponding representation-action, which produces just that representation which
answers to the object. The individual organism, which contains within itself the
real possibility of all possible representation-productions, is as it were
calculated to work, at every point---that is, in every respect---in such accord or
concordance with every influence coming from the side of the total, that the
former each time produces precisely what the latter aims to call forth (as when an
instrument, according to its own entire construction, renders---though in its own
individual way---precisely the tone that is struck on another). We feel ourselves,
however, in this not as related to mere counter-images, but as related to actual
objects, because we here constantly feel our organism subjected to an affection
which, coming from something other than ourselves, is what calls forth and compels
that action and production of our own soul---so that the affection of our
individual organism, or microcosm, distinguishes for us sensation from mere
representation.\par}

\begin{center}\large\S~4.\end{center}
\markboth{\S~4}{\S~4}
\begin{center}Philosophy is b) intelligent and rational cognizing. ---The fundamental
moments of cognition in their appearance as such. ---Cognition as intelligent.\end{center}

The individual life of consciousness, determined by the power of the objective,
shows itself, insofar as it comes forward over against the objective, as that
which is in itself free, carrying with it an inexhaustible possibility of acting
in other ways than the way which the determining necessity each time brings with
it. The necessity thus refers here to a freedom, and has such a freedom for its
object, since it is precisely the in-itself-free (the in-itself-free imagination
and reflection) that is bound by the necessity; and only thereby does cognition
acquire the character of a cognizing, since only thereby does it show itself as an
ascription of validity [Gyldighedstillæggelse]. [The true, too, comes forward as
the valid and as that which is cognized in its validity only to the same degree as
this free element, coming forward as such, then shows itself to be that which is
determined and bound.]

But in the first natural cognizing, which forms itself of its own accord, and in
which we on the whole comport ourselves as natural beings (cf.\ my
\textit{Psych., ny Udarb.}, \S~26), the opposition between the objective and the
subjectivity, though it may show itself in particular cases, does not come forward
in such a way that it wholly asserts itself as such. Cognition finds itself then,
despite all the mediating element lying within it, at a first stage of immediacy,
and the moments of cognition, though certainly present, are so lost in the
cognizing as a whole that they do not come forward as such. But then neither does
the true come forward for us in \emph{its truth}, nor does it come to assert
itself by its formally acknowledged validity.

In order for the true to attain such a higher governance, the individual life of
consciousness lying on the side of the subjectivity must first come wholly forth
into its freedom. And to this, too, the striving for cognition leads in its
advance, in that we---whether by some other interest, or by mere desire to know
[Videdrivt] and the higher life working itself forth through it---are moved to go
beyond the immediately given, to seek other and more: either merely a) by looking
further about us in the sphere opened to us, whereby not only the directing of the
mind, that is, attention, becomes free, but also, as expectation sets in, the
representation begins to find itself unbound; or else b) by asking after other and
more than what is given in that sphere, because through it we find ourselves
awakened to seek something such as that to which the given points---so that a
spirit of inquiry [Randsagelsesaand] asserts itself, and scrutiny [Granskning] or
investigation [Efterforskning] comes forth, and now not only the representation,
the imagination, but also the pursuit of relations and connections lying within
it, or reflection, comes forward as such. ---But the new that is now in this way
formally sought is sought at the same time as the valid, and as answering to the
investigation in such a way that that questioning can find itself satisfied;
whereby one strives after a result by which the representing and pursuing that has
become free---or has stepped forth into its freedom---can again be sufficiently
determined and bound, with a necessity that is now also sought, and which, when
found, will come to come forward as such.

But the necessity---together with the being-governed by, and being-at-home in, the
power of the objective---that sets in in this way has then a different and freer
character than that first immediate necessity. It now comes forward as necessity
of understanding and reason [Forstands- og Fornuftnødvendighed], and we ourselves
now show ourselves as beings of reason who ascribe validity because we acknowledge
according to an express admission, assent, and approval grounded in insight;
whereby the acknowledgment, or ascription of validity, too comes---stepping forth
into its freedom---to show itself as such. For although at first it may be only new
objects, of the same kind as those given in an immediate way, that are thus sought
under a merely tracking observation, they must now not only show themselves to us
as the ones we sought, but we must now also be assured that they are the right
ones, and that we now objectify them to ourselves rightly, so that nothing else
presents itself to us and is posited by us than the in-itself-valid, or that which
\emph{ought} to be posited. ---But this then naturally has such a retroactive
effect on what was earlier immediately assumed that it too can be drawn into the
same consideration, so that \emph{all} cognizing can become the object of a
question as to its validity.

Cognition will in this way become one that we ourselves confirm in ourselves, in
that we confirm ourselves in it, so that now the assurance too comes to show itself
as such, and indeed as a potentiated assurance, or as an assurance of whose
validity we are assured. ---But certainly the uncertainty lying in the questioning
may here too come forward in such a way that something doubting, dubitative, may
set in, which is however overcome when a result is attained that becomes
sufficiently satisfying.

Here, then, there emerges a properly mediated cognizing---mediated, namely, by
thinking, and resting on thinking and the satisfaction of thinking---an intelligent
cognizing, in the formation of which the fundamental demands originally lying in
the pursuit of connection and relation that strives toward cognition, that is, in
thinking itself (in the intelligence in the subject), are so in motion that it is
in accordance with their satisfaction that insight, approval, and acknowledgment
set in.

\begin{center}\large\S~5.\end{center}
\markboth{\S~5}{\S~5}
% [text to be added: pp. 15--20 (§ 5): Fortsættelse. Hvorpaa den intelligente Erkjenden gaaer ud.]

\begin{center}\large\S~6.\end{center}
\markboth{\S~6}{\S~6}
% [text to be added: p. 21 (§ 6): Erkjendelsen som rational, henarbeidende til et sidste Grundlag.]

\begin{center}\large\S~7.\end{center}
\markboth{\S~7}{\S~7}
% [text to be added: pp. 22--24 (§ 7): Philosophie er c) videnskabelig Erkjendelse. Om Videnskab og Videnskabelighed i Almindelighed.]

\begin{center}\rule{0.18\linewidth}{0.4pt}\end{center}

%% ============================================================
%%  THE PROPAEDEUTIC CONCEPT OF PHILOSOPHY AND ITS EXPOSITION. (§§ 8--14, pp. 25--58)
%% ============================================================
\grouphead{The Propaedeutic Concept of Philosophy and Its Exposition.}

\begin{center}\large\S~8.\end{center}
\markboth{\S~8}{\S~8}
% [text to be added: pp. 25--26 (§ 8): Philosophie gaaer d) ud paa en altomfattende Grunderkjendelse. Subjectiv Hovedopgave.]

\begin{center}\large\S~9.\end{center}
\markboth{\S~9}{\S~9}
% [text to be added: p. 27 (§ 9): Objectiv Hovedopgave. --- Philosophie gaaer e) ud paa Tilværelsens sidste Grunde.]

\begin{center}\large\S~10.\end{center}
\markboth{\S~10}{\S~10}
% [text to be added: pp. 28--32 (§ 10): Propædeutisk Definition paa Philosophie.]

\begin{center}\large\S~11.\end{center}
\markboth{\S~11}{\S~11}
% [text to be added: pp. 33--41 (§ 11): Philosophien træder det ellers Givne imøde. --- Overblik over Indbegrebet af alt dette.]

\begin{center}\large\S~12.\end{center}
\markboth{\S~12}{\S~12}
% [text to be added: p. 42 (§ 12): Exposition af det givne Begreb om Philosophie: tre Hovedsynspuncter herved.]

\begin{center}\large\S~13, a.\end{center}
\markboth{\S~13a}{\S~13a}
% [text to be added: pp. 42--46 (§ 13, a): Exposition af Begrebet Philosophie i første Hovedhenseende.]

\begin{center}\large\S~13, b.\end{center}
\markboth{\S~13b}{\S~13b}
% [text to be added: pp. 47--56 (§ 13, b): Begrebets Exposition i anden Hovedhenseende.]

\begin{center}\large\S~14.\end{center}
\markboth{\S~14}{\S~14}
% [text to be added: pp. 56--58 (§ 14): Begrebets Exposition i tredie Hovedhenseende.]

%% ============================================================
%%  THE FOUNDATION OF PHILOSOPHY AND ITS WORK. (§§ 15--19, pp. 59--80)
%% ============================================================
\grouphead{The Foundation of Philosophy and Its Work.}

\begin{center}\large\S~15.\end{center}
\markboth{\S~15}{\S~15}
% [text to be added: pp. 59--62 (§ 15): Den philosophiske Grundidee.]

\begin{center}\large\S~16.\end{center}
\markboth{\S~16}{\S~16}
% [text to be added: pp. 63--66 (§ 16): Hvori den philosophiske Grundidee har sit Grundlag.]

\begin{center}\large\S~17.\end{center}
\markboth{\S~17}{\S~17}
% [text to be added: pp. 66--67 (§ 17): Hvorledes Philosophie kommer i Stand. --- a) Explication i speculativ Tendents.]

\begin{center}\large\S~18.\end{center}
\markboth{\S~18}{\S~18}
% [text to be added: pp. 67--71 (§ 18): b) Speculation og Dialektik i væsentlig Eenhed.]

\begin{center}\large\S~19.\end{center}
\markboth{\S~19}{\S~19}
% [text to be added: pp. 71--80 (§ 19): Philosophiens Udtræden i c) et philosophiskt System.]

%% ============================================================
%%  THE FULL DETERMINATION OF THE CONCEPT OF PHILOSOPHY. (§ 20, p. 81)
%% ============================================================
\grouphead{The Full Determination of the Concept of Philosophy.}

\begin{center}\large\S~20.\end{center}
\markboth{\S~20}{\S~20}
% [text to be added: p. 81 (§ 20): Endelig Definition paa Philosophie.]

%% ============================================================
%%  PHILOSOPHY AS A MEMBER AND MOMENT IN THE WHOLE SPIRITUAL LIFE. (§ 21, pp. 82--86)
%% ============================================================
\grouphead{Philosophy as a Member and Moment in the Whole Spiritual Life.}

\begin{center}\large\S~21.\end{center}
\markboth{\S~21}{\S~21}
% [text to be added: pp. 82--86 (§ 21): Philosophie paa lige Linie med det Empiriske og det Suprarationale, --- med Poesie, Moral og Religion.]

%% ============================================================
%%  RETTELSER (errata leaf). The Danish errata correct only the ORIGINAL's
%%  orthography, so they are meaningless in English. Per the erkjendelse policy,
%%  render a one-line translator's note here at finish (or omit). Decide at finish.
%% ============================================================
% [text to be added: Rettelser -- translator's note (see TRANSLATION-NOTES finish checklist)]

\end{document}
